\documentclass[sigconf,nonacm]{acmart}
\settopmatter{printacmref=false}
\renewcommand\footnotetextcopyrightpermission[1]{}
\pagestyle{plain}

\begin{document}

\title{Expense Tracker: A Native Android App with Jetpack Compose, Room, and Intents}

\author{Student Name 1}
\affiliation{\institution{Ain Shams University}\city{Cairo}\country{Egypt}}
\email{student1@example.edu}

\author{Student Name 2}
\affiliation{\institution{Ain Shams University}\city{Cairo}\country{Egypt}}
\email{student2@example.edu}

\author{Student Name 3}
\affiliation{\institution{Ain Shams University}\city{Cairo}\country{Egypt}}
\email{student3@example.edu}

\begin{abstract}
Many students and young professionals lose track of small daily purchases, which makes monthly budgeting difficult. This project presents \emph{Expense Tracker}, a native Android application, written in Kotlin, that lets users record expenses, organize them by category, and see a running total and a per-category summary at a glance. The interface is built entirely with Jetpack Compose and follows Material Design 3. Expenses are stored locally with the Room persistence library, and the user's preferred currency is stored with Preferences DataStore, so the app works offline and keeps its data between sessions. The app also communicates with other applications through Android Intents: an explicit Intent opens a details screen for a selected expense, while implicit Intents let the user share a spending summary and look up a nearby ATM in a maps application. This abstract describes the problem, the architecture, the main Android concepts used, and the planned evaluation.
\end{abstract}

\keywords{Android, Kotlin, Jetpack Compose, Room, DataStore, Intents, personal finance}

\maketitle

\section{Introduction}
Personal expense tracking is a common use case for mobile computing because a phone is always at hand when a purchase happens. Existing finance apps are often feature-heavy, require accounts, or depend on a network connection. Our goal is a small, offline-first app that covers the essential workflow: add an expense, find it later, correct it, and understand where the money goes. The project also serves as a practical exercise in the core concepts taught in the Mobile Computing course: the Activity lifecycle, declarative UI, local persistence, and inter-app communication.

\section{Features}
The application supports the following functions.
\begin{itemize}
  \item \textbf{CRUD operations.} Users can create, view, edit, and delete expenses. Each expense has a title, an amount, a category (Food, Transport, Bills, Shopping, Health, Other), a date, and an optional note. Input is validated so that an empty title or a non-positive amount cannot be saved.
  \item \textbf{Summary.} The home screen shows the total spent and the top categories, and updates automatically whenever the data changes.
  \item \textbf{Currency preference.} The user can switch between EGP, USD, EUR, and SAR; the choice is remembered across launches.
  \item \textbf{Sharing and discovery.} The user can share a text summary with any compatible app and open a maps application to find a nearby ATM.
\end{itemize}

\section{Architecture and Android Concepts}
The app follows a simple layered design with unidirectional data flow, shown in Table~\ref{tab:layers}.

\begin{table}[h]
  \caption{Layers of the application.}
  \label{tab:layers}
  \begin{tabular}{p{1.9cm}p{5.6cm}}
    \toprule
    Layer & Responsibility \\
    \midrule
    UI & Composable screens and a reusable add/edit dialog; they only observe state and send events. \\
    ViewModel & Holds UI state as \texttt{StateFlow} and launches database writes in \texttt{viewModelScope}. \\
    Data & Room entity, DAO, and database; DataStore for preferences. \\
    \bottomrule
  \end{tabular}
\end{table}

\paragraph{Jetpack Compose.} Screens are written as composable functions. State is hoisted to the ViewModel and collected with lifecycle awareness, so the UI is recomputed only when the underlying data changes and survives configuration changes such as rotation.

\paragraph{Persistence.} Room exposes expenses as a \texttt{Flow}, so every insert, update, or delete is reflected in the UI without manual refreshing. A single \texttt{Upsert} DAO method covers both creation and editing. Preferences DataStore keeps the currency setting, which suits small key-value data better than a relational table.

\paragraph{Intents.} Navigation from the list to the details screen uses an \emph{explicit} Intent that carries the expense identifier as an extra. Sharing and the ATM lookup use \emph{implicit} Intents (\texttt{ACTION\_SEND} and \texttt{ACTION\_VIEW} with a \texttt{geo:} URI). Because the app cannot know which apps are installed, each implicit call handles the case where no application responds.

\paragraph{Code quality.} The code uses immutable \texttt{val} declarations and data classes where possible, extension functions, and lambdas for event handling, in line with idiomatic Kotlin.

\section{Planned Evaluation}
We will test the main flows manually on an emulator and on a physical device: adding, editing, and deleting expenses; rotating the device; restarting the app to confirm that data and the currency choice persist; and triggering both kinds of Intents. We will also gather informal feedback from classmates on the clarity of the interface. Possible extensions for the next phase include monthly filtering, charts, and data export.

\section{Conclusion}
Expense Tracker demonstrates how a small but complete Android application can combine declarative UI, local persistence, and Intents in a clean, maintainable structure. It gives users a fast way to record spending while giving the team hands-on experience with the platform concepts covered in the course.

\end{document}
